\chapter{\IfLanguageName{dutch}{Stand van zaken}{State of the art}}
\label{ch:stand-van-zaken}

Dit hoofdstuk bespreekt de literatuur die relevant is voor het optimaliseren van de initiële projectopstart van frontend-applicaties binnen een gestandaardiseerde agency-omgeving. Waar de inleiding de praktische uitdagingen rond repetitief werk en inconsistentie introduceerde, verkent dit hoofdstuk mogelijke aanpakken om dit proces te verbeteren. Eerst wordt de brede context van AI-ondersteunde codegeneratie en het daadwerkelijke gebruik in de praktijk geschetst. Vervolgens wordt dieper ingegaan op de integratie van domeinspecifieke context via standaarden zoals het Model Context Protocol (MCP) en het gebruik van \textit{agentic workflows}. Tot slot belicht dit hoofdstuk de impact van deze aanpakken op technische schuld en de noodzaak van objectieve evaluatiemetrieken. Deze literatuur vormt de theoretische basis om te bepalen onder welke voorwaarden een context-gedreven aanpak de beste balans biedt tussen efficiëntie, kwaliteit en consistentie.

\section{LLM-codegeneratie in software engineering}
\label{sec:llm-codegeneratie}

Recente overzichtsstudies zoals \textit{A Survey on Large Language Models for Code Generation} tonen aan dat Large Language Models (LLM's) inmiddels breed worden ingezet voor diverse softwaretaken, gaande van algemene codegeneratie tot testautomatisering. Hoewel deze modellen over de hele lijn efficiëntiewinst opleveren, blijkt uit de \textit{Systematic Survey on Large Language Models for Code Generation} dat de betrouwbaarheid, effectiviteit en beveiliging van de gegenereerde code grote uitdagingen blijven. Deze beperkingen maken dat algemene modellen vaak moeite hebben met het afleveren van direct inzetbare output zonder manuele correcties. Voor dit onderzoek betekent dit dat de focus niet louter op de generatie van code kan liggen, maar dat de betrouwbaarheid binnen een specifieke projectcontext primair staat.

Daarnaast wijst de \textit{Survey on LLM-based Code Generation for Low-Resource and Domain-Specific Programming Languages} erop dat het gebruik van LLM's in domeinspecifieke omgevingen met beperkte trainingsdata bijzonder complex is. Algemene modellen missen vaak de nodige context over interne componentenbibliotheken en architecturale afspraken, en standaard benchmarks ontbreken om deze specifieke prestaties te evalueren. Dit is cruciaal voor de probleemstelling van Ballistix Digital, waar de projectopstart sterk afhankelijk is van een gesloten, interne React-bibliotheek. Het benadrukt de noodzaak om te onderzoeken in welke mate extra context de kwaliteit en bruikbaarheid van de output verhoogt.

\section{Agentic workflows en tool-use}
\label{sec:agentic-workflows}

Naast het aanleveren van passieve context, verschuift de literatuur naar actieve, agent-gedreven benaderingen. De paper \textit{From LLMs to LLM-based Agents for Software Engineering} stelt dat het toevoegen van autonomie, besluitvormingscapaciteiten en tool-integratie een taalmodel transformeert in een AI-agent. Deze agents kunnen complexe taken, zoals de opstart van een volledig project of testgeneratie, opdelen in kleinere stappen en iteratief uitvoeren. Deze autonomie introduceert echter ook onvoorspelbaarheid; een agent kan foute beslissingen nemen of in een eindeloze lus terechtkomen wanneer de feedback loop ontbreekt. Dit onderstreept het belang voor dit onderzoek om te evalueren onder welke voorwaarden \textit{agentic workflows} daadwerkelijk tijdswinst opleveren bij de projectopstart zonder in te boeten op stabiliteit.

Daarnaast beargumenteert de studie \textit{If LLM Is the Wizard, Then Code Is the Wand} dat de integratie van code en tool-gebruik cruciaal is voor het redeneer- en planningsvermogen van LLM's. Modellen die acties kunnen uitvoeren (zoals bestanden lezen of commando's draaien), blijken beter in staat om complexe engineeringstaken af te ronden dan modellen die enkel tekst genereren. Een kanttekening hierbij is dat dit ook hogere eisen stelt aan de kwaliteitscontrole van de uitgevoerde acties. Voor de context van dit onderzoek toont dit aan dat het gebruik van tools essentieel is om de kloof tussen theorie en werkende projectcode te overbruggen.

Tot slot is het belangrijk om te kijken naar hoe ontwikkelaars deze systemen in de praktijk gebruiken. Uit \textit{Beyond Code Generation: An Observational Study of ChatGPT Usage in Software Engineering Practice} blijkt dat ontwikkelaars LLM's vaker gebruiken als hulpmiddel voor uitleg, brainstorming of kleine iteraties, in plaats van als volledige vervanging voor het schrijven van complexe, samenhangende codebases. Dit toont aan dat een volledig geautomatiseerde projectopstart wellicht minder realistisch is dan een aanpak waarbij de ontwikkelaar ondersteund wordt. Voor dit onderzoek rechtvaardigt het de keuze voor een iteratieve, context-gedreven aanpak waarbij de ontwikkelaar de regie behoudt.

\section{Context-integratie en het Model Context Protocol (MCP)}
\label{sec:mcp}

Om de beperkingen van generieke modellen in domeinspecifieke omgevingen te overbruggen, wordt toenemend gezocht naar manieren om context en tools te structureren. De \textit{Model Context Protocol specification} beschrijft een open standaard op basis van een host-client-server model die precies dit doel dient. In plaats van ad-hoc oplossingen te bouwen voor elke databron, maakt MCP het mogelijk om op een gestandaardiseerde manier bedrijfsinterne context, \textit{Resources}, \textit{Tools} en \textit{Prompts} te delen, aldus de \textit{MCP documentation}. Hierdoor kunnen modellen redeneren over en handelen binnen een gecontroleerde omgeving. Toch moet worden opgemerkt dat MCP relatief nieuw is en niet op zichzelf de bedrijfslogica oplost. Voor dit onderzoek vormt de standaard echter een zeer relevante oplossingsrichting om te testen in welke mate een gestructureerde ontsluiting van de interne componentenbibliotheek en architectuur de efficiëntie van de projectopstart verbetert ten opzichte van traditionele prompts.

\section{Technische schuld en onderhoudbaarheid}
\label{sec:technische-schuld}

Een kernvraag bij het inzetten van generatieve AI is de langetermijnimpact op de codebase. Publicaties zoals \textit{How AI generated code compounds technical debt} wijzen erop dat AI-output de initiële ontwikkelsnelheid kan verhogen, maar tegelijk de onderhoudslast aanzienlijk kan vergroten als de gegenereerde code afwijkt van architecturale standaarden. Ook bronnen zoals \textit{Why AI-generated code is creating a technical debt nightmare} en \textit{The Hidden Costs of Coding With Generative AI} benadrukken dat de kortetermijnwinst in productiviteit vaak ten koste gaat van langetermijncomplexiteit. Deze literatuur waarschuwt dat code die oppervlakkig correct lijkt, structureel zwak kan zijn of niet herbruikbaar is binnen de bestaande context. Voor de agency-context van dit onderzoek is dit inzicht cruciaal: een versnelde projectopstart heeft enkel waarde als het resultaat geen extra correctiewerk oplevert of de consistentie van de codebase in gevaar brengt. Het valideert bovendien de keuze om niet alleen de opleversnelheid te meten, maar expliciet te sturen op kwaliteitsbehoud.

\section{Evaluatiemetrieken voor codekwaliteit}
\label{sec:evaluatiemetrieken}

Om de impact op technische schuld en efficiëntie kwantificeerbaar te maken, is een duidelijke evaluatiemethode noodzakelijk. Uit voornoemde studies naar LLM-codegeneratie blijkt dat veelgebruikte metrieken zoals Pass@k en BLEU-scores weliswaar de functionele correctheid of tekstuele gelijkenis meten, maar onvoldoende de structurele kwaliteit of de mate van technische schuld vatten. Dit betekent dat traditionele AI-benchmarks tekortschieten voor het evalueren van een projectopstart, waarbij architecturale integriteit en leesbaarheid vooropstaan. Voor dit onderzoek is het daarom noodzakelijk om een bredere set van metrieken te hanteren die de balans tussen snelheid en kwaliteit weerspiegelen, zoals het meten van het benodigde correctiewerk, de totale doorlooptijd en de mate van consistentie met de huisstijl van de agency.